\section{三个实证锚点}

这些索引来源在实证上的作用是狭义且有纪律的：
\begin{enumerate}[nosep]
  \item \textbf{意识状态的标志。} 2024 年的 Communications
        Biology 文献只用于支撑一个强主张：意识条件与稳健且可泛化的动力学结构--功能标志相关，
        而不只是能够用松散的心理学语言来描述。
  \item \textbf{连接性与稳定性。} 2025 年的 Communications Biology
        文献用于支撑第二个强主张：全局意识状态与有结构的连接性、局部动力学和稳定性有关，
        因而状态持续并非单纯噪声问题，而是一个架构问题。
  \item \textbf{转变动力学。} 2025 年的 Nature Neuroscience 睡眠文献
        用于支撑第三个强主张：进入一个新机制可以是突变式的、有模式可循的，
        并且能够用动力学语言处理，而不必停留在纯描述层面。
\end{enumerate}

\begin{discussion}
这些来源并不能证明完整的框架惯性模型，本章也不应假装如此。
它们的价值更为克制：它们为把状态、稳定性与转变视为严肃的生物学议题提供了正当性。
一旦这个底板成立，手稿其余部分再去讨论进入成本、机制持续与重定向，
就不会听起来像是脱离有组织状态变化科学的私人生产力隐喻。
\end{discussion}